% Nouveau EGA V (cote 25): archive cote 25, site page 91 (PDF page 92).
% typescript. Draft compared: none.
% AI-transcribed by Claude Opus 5.5 (claude-opus-5-5, adaptive thinking, effort high), Message Batches API, batch ega25b (msgbatch_01A7n5YuXubUPNK1zsepi8RG), 2026-11-13; not yet checked by eye.
\scanpage{25/91}
\mspage{85}

Bien entendu, pour justifier dans le cas \xout{\unsure{actuel}} du présent n° le passage
de \hand{$\eta$} générique aux points voisins de $\mathrm{Grass}^{m+1}(P)$, et pouvoir également
le cas échéant reprendre à notre compte les explications générales de
7.1. , il y a lieu de considérer les diagrammes
\diagram{\begin{tikzcd}
X \arrow[d] & \widetilde{X}(C) \arrow[l] \arrow[d] \\
S & Q(C) \arrow[l]
\end{tikzcd}}
obtenus à l'aide des divers $C \in \mathrm{Grass}^{m+1}(S)$, et \inserted{plus généralement} ceux obtenus de même
après changement de base $T \to S$ pour des points \hand{$\zeta$} \hand{$\in$} $\mathrm{Grass}^{m+1}($\over{P}{\unsure{T}}$)$ :
\diagram{\begin{tikzcd}
X_T \arrow[d] & \widetilde{X}(C_{\hand{$\zeta$}}) = \widetilde{X_T}(C_{\hand{$\zeta$}}) \arrow[l] \arrow[d] \\
T & Q(C_{\hand{$\zeta$}}) \arrow[l]
\end{tikzcd}}
comme déduits, par changement\hand{s} de base \hand{$\zeta$} : $T \longrightarrow \mathrm{Grass}^{m+1}$\deleted{$(P)$} $= \underline{T}$ ,
du diagramme \uline{universel} (relatif au point canonique de $\mathrm{Grass}^{m+1}$ dans $\underline{T}$ ):
\diagram{\begin{tikzcd}
X_{\underline{T}} \arrow[d] & \widetilde{X}(\underline{C}) \arrow[l] \arrow[d] \\
\underline{T} & Q(\underline{C}) \arrow[l]
\end{tikzcd}}
où $\underline{C}$ est la sous-variété linéaire canonique de $P_{\underline{T}}$ . Alors ci-dessus
$\widetilde{X}(C_{\hand{\unsure{$\eta$|$\zeta$}}}) \longrightarrow Q(C_{\hand{\unsure{$\eta$|$\zeta$}}})$ n'est autre que le morphisme des fibres génériques pour
le $\underline{T}$-morphisme $\widetilde{X}(\underline{C}) \longrightarrow Q(\underline{C})$ de ce dernier diagramme, et toute propri\'e-%
té constructible pour \xout{xxxx} le morphisme des fibres génériques implique
la propriété pour les fibres voisines. - Du point de vue notationel,
$Q$ doit être envisagé (et même introduit) comme le nom du morphisme \inserted{de foncteurs}
naturel $\mathrm{Grass}^{m+1}(P) \longrightarrow$ \xout{Hilb} Sous-préschémas $(\mathrm{Grass}^{m}(P))$ .

15. \uline{Axiomatisation de certains des résultats précédents}.

Je pense surtout aux résultats des n°s 2 à 8, qui pour la plupart

\mspage{91}
